% Created 2026-06-01 Mon 23:55
% Intended LaTeX compiler: pdflatex
\documentclass[11pt, letterpaper]{article}
\usepackage[utf8]{inputenc}
\usepackage[T1]{fontenc}
\usepackage{graphicx}
\usepackage{longtable}
\usepackage{wrapfig}
\usepackage{rotating}
\usepackage[normalem]{ulem}
\usepackage{amsmath}
\usepackage{amssymb}
\usepackage{capt-of}
\usepackage{hyperref}
\usepackage[margin=2.4cm]{geometry}
\usepackage[utf8]{inputenc}
\usepackage[T1]{fontenc}
\usepackage{lmodern}
\usepackage{xcolor}
\usepackage{hyperref}
\hypersetup{colorlinks=true, linkcolor=blue, urlcolor=blue}
\author{Equipo Alpine White}
\date{\textit{{[}2026-04-21 Tue]}}
\title{Bitácora Semanal 11: Administración de Sistemas Unix/Linux}
\hypersetup{
 pdfauthor={Equipo Alpine White},
 pdftitle={Bitácora Semanal 11: Administración de Sistemas Unix/Linux},
 pdfkeywords={},
 pdfsubject={},
 pdfcreator={Emacs 30.2 (Org mode 9.7.11)}, 
 pdflang={English}}
\begin{document}

\maketitle
\setcounter{tocdepth}{2}
\tableofcontents

\section{Información General}
\label{sec:orga75b20f}

\begin{itemize}
\item \textbf{Semana:} Semana 11
\item \textbf{Materia:} Administración de Sistemas Unix/Linux
\item \textbf{Docente:} Luis Enrique Serrano Gutiérrez
\item \textbf{Equipo:}
\begin{itemize}
\item Arreguín Salgado Gael Emiliano
\item López Pérez Mariana
\item Nieto Gallegos Isaac Julián
\end{itemize}
\end{itemize}

Durante esta semana se trabajó con tres tecnologías importantes dentro de la administración de sistemas Linux: sincronización de tiempo mediante NTP y Chrony, montaje automático de recursos con AutoFS y administración flexible de almacenamiento mediante LVM.

El objetivo principal fue comprender cómo Linux automatiza tareas críticas mediante servicios, archivos de configuración y capas de abstracción. Estas herramientas permiten que el sistema sea más estable, escalable y eficiente en ambientes de servidor.
\subsection{Responsabilidades}
\label{sec:orgaa071f7}
\begin{itemize}
\item \textbf{Gael Emiliano:} Documentación de configuración de Chrony, comandos NTP y explicación del servidor de hora.
\item \textbf{Mariana:} Captura de conceptos vistos en clase y ampliación de AutoFS, NFS y recursos remotos.
\item \textbf{Isaac Julián:} Revisión del archivo Org, bloques de código, formato y compilación final.
\end{itemize}

\newpage
\section{Sincronización de tiempo con NTP y Chrony}
\label{sec:org45a6d27}

\subsection{Conceptos nuevos}
\label{sec:orga5ae4cf}

\subsubsection{Servidor}
\label{sec:org0380202}

Un servidor es un sistema físico o lógico que proporciona servicios, datos o recursos a otros equipos llamados clientes.

En esta práctica se trabajó con un servidor NTP, encargado de permitir que otros equipos de la red sincronicen su reloj con una fuente de tiempo confiable.
\subsubsection{NTP}
\label{sec:org0ee3450}

NTP significa \emph{Network Time Protocol}. Es un protocolo utilizado para sincronizar relojes de computadoras, servidores y dispositivos de red.

Mantener la hora correcta es importante porque afecta:

\begin{itemize}
\item Logs del sistema.
\item Auditorías.
\item Certificados digitales.
\item Autenticación.
\item Bases de datos.
\item Tareas programadas.
\item Servicios distribuidos.
\end{itemize}

Sin una hora sincronizada, los eventos del sistema pueden quedar registrados en orden incorrecto, dificultando diagnósticos y revisiones de seguridad.
\subsubsection{Chrony}
\label{sec:org889a491}

Chrony es una implementación moderna del protocolo NTP.

Está formada principalmente por:

\begin{itemize}
\item \texttt{chronyd}: demonio encargado de sincronizar el reloj.
\item \texttt{chronyc}: herramienta para consultar y administrar Chrony.
\item \texttt{/etc/chrony.conf}: archivo principal de configuración.
\end{itemize}

Chrony puede funcionar como cliente NTP o como servidor NTP para otros equipos dentro de una red local.
\subsubsection{chronyd}
\label{sec:org70bd664}

\texttt{chronyd} es el servicio que se ejecuta en segundo plano y ajusta el reloj del sistema.

\begin{verbatim}
systemctl status chronyd
\end{verbatim}

Después de modificar la configuración, se debe reiniciar:

\begin{verbatim}
sudo systemctl restart chronyd
\end{verbatim}
\subsubsection{chronyc}
\label{sec:org326c24b}

\texttt{chronyc} permite consultar el estado de Chrony.

\begin{verbatim}
chronyc sources -v
\end{verbatim}

Este comando muestra información sobre:

\begin{itemize}
\item Fuente de sincronización.
\item Stratum.
\item Retardo.
\item Precisión.
\item Estado de cada servidor.
\end{itemize}
\subsubsection{Stratum}
\label{sec:orgf5dd23a}

El stratum indica la distancia jerárquica respecto a una fuente original de tiempo.

\begin{itemize}
\item \textbf{Stratum 0:} reloj atómico o GPS.
\item \textbf{Stratum 1:} servidor conectado directamente al Stratum 0.
\item \textbf{Stratum 2:} servidor sincronizado con Stratum 1.
\item \textbf{Stratum 3:} servidor sincronizado con Stratum 2.
\end{itemize}

Mientras menor sea el stratum, más cerca se encuentra la fuente original de tiempo.
\subsection{Configuración de Chrony como servidor NTP}
\label{sec:org4017183}

El archivo principal de configuración es:

\begin{verbatim}
/etc/chrony.conf
\end{verbatim}

Ejemplo de configuración:

\begin{verbatim}
server ntp.lab.int iburst
allow 192.168.0.0/24
driftfile /var/lib/chrony/drift
makestep 1.0 3
rtcsync
logdir /var/log/chrony
\end{verbatim}

La directiva \texttt{allow} permite que los clientes de una subred consulten al servidor NTP.

Después de configurar Chrony:

\begin{verbatim}
sudo systemctl restart chronyd
\end{verbatim}

También se debe permitir NTP en el firewall:

\begin{verbatim}
sudo firewall-cmd --add-service=ntp --permanent
sudo firewall-cmd --reload
\end{verbatim}
\subsection{Conceptos de repaso}
\label{sec:orge6d4f76}

\begin{itemize}
\item Un servicio en Linux suele ejecutarse como demonio.
\item \texttt{systemctl} permite administrar servicios.
\item Los archivos en \texttt{/etc} contienen configuración persistente.
\item Un servidor puede compartir tiempo, archivos, datos o autenticación.
\item La sincronización de tiempo facilita auditorías y diagnóstico.
\end{itemize}
\subsection{Máximas}
\label{sec:org4c734a8}
\begin{itemize}
\item “Chrony no solo sincroniza el reloj; elige dinámicamente la fuente más precisa disponible en la red.”
\item “Un servidor con hora incorrecta registra eventos correctos en momentos incorrectos.”
\item “El stratum indica cercanía a la fuente original, pero la calidad real también depende de la red.”
\item “NTP no es decoración del sistema; es base para auditoría, seguridad y diagnóstico.”
\end{itemize}
\subsection{Sección libre}
\label{sec:org81bb5b9}

\texttt{ntpd} fue durante mucho tiempo el servicio tradicional de sincronización NTP, pero en sistemas modernos Chrony suele ser preferido por su rapidez y precisión.

\texttt{ntpdate} permite sincronización puntual, aunque actualmente está en desuso.

También puede revisarse el estado general del tiempo con:

\begin{verbatim}
timedatectl
\end{verbatim}

\newpage
\section{AutoFS y montaje automático de recursos}
\label{sec:org1400f5a}

\subsection{Conceptos nuevos}
\label{sec:org0bc143d}

\subsubsection{AutoFS}
\label{sec:orgde84e3a}

AutoFS es un servicio que permite montar sistemas de archivos automáticamente bajo demanda.

A diferencia de \texttt{/etc/fstab}, AutoFS no monta recursos al iniciar el sistema. Solo los monta cuando un usuario o proceso intenta acceder a ellos.

Esto es útil para recursos remotos porque evita que el arranque dependa de servidores externos.
\subsubsection{Diferencia entre fstab y AutoFS}
\label{sec:org6b72024}

\texttt{/etc/fstab} monta recursos durante el arranque.

AutoFS funciona de otra manera:

\begin{itemize}
\item Espera a que se acceda al recurso.
\item Monta automáticamente el directorio.
\item Mantiene el recurso disponible mientras se usa.
\item Lo desmonta después de un tiempo de inactividad.
\end{itemize}

Esta lógica reduce consumo de recursos y evita bloqueos si el servidor remoto no está disponible.
\subsubsection{auto.master}
\label{sec:org6190d3b}

\texttt{/etc/auto.master} es el archivo principal de AutoFS.

Ejemplo:

\begin{verbatim}
/home/remoto /etc/auto.users
\end{verbatim}

Esto indica que los montajes dentro de \texttt{/home/remoto} serán definidos por el archivo \texttt{/etc/auto.users}.
\subsubsection{Mapas secundarios}
\label{sec:orgf43fb84}

Los mapas secundarios definen qué recurso remoto se monta.

Ejemplo:

\begin{verbatim}
usuario -rw servidor:/export/home/usuario
\end{verbatim}

Con esta regla, AutoFS monta el recurso cuando se accede al directorio correspondiente.
\subsubsection{Timeout}
\label{sec:org72fb13d}

El timeout permite desmontar recursos inactivos automáticamente.

Esto no representa un error, sino una forma de liberar recursos de red y evitar conexiones innecesarias.
\subsection{Conceptos de repaso}
\label{sec:org54083b8}

\begin{itemize}
\item NFS permite compartir directorios en red.
\item \texttt{mount} realiza montajes manuales.
\item \texttt{/etc/fstab} realiza montajes persistentes.
\item AutoFS monta bajo demanda.
\item El montaje integra recursos externos al árbol local de directorios.
\end{itemize}
\subsection{Máximas}
\label{sec:org6c503f6}

\begin{itemize}
\item “AutoFS no monta cuando arranca el sistema; monta cuando alguien realmente necesita el recurso.”
\item “Un timeout en AutoFS no es un error; es eficiencia.”
\item “Montar bajo demanda evita depender de recursos remotos durante el arranque.”
\item “AutoFS convierte recursos de red en accesos casi transparentes para el usuario.”
\end{itemize}
\subsection{Sección libre}
\label{sec:orgdc7bce0}

Para reiniciar AutoFS:

\begin{verbatim}
sudo systemctl restart autofs
\end{verbatim}

Para revisar su estado:

\begin{verbatim}
systemctl status autofs
\end{verbatim}

AutoFS es útil en laboratorios, servidores multiusuario y redes donde se comparten directorios mediante NFS.

\newpage
\section{LVM y administración dinámica de almacenamiento}
\label{sec:org29ef1a5}

\subsection{Conceptos nuevos}
\label{sec:org3c2b420}

\subsubsection{LVM}
\label{sec:orge234594}

LVM permite administrar almacenamiento mediante capas lógicas en lugar de depender directamente de particiones físicas rígidas.

Su estructura principal es:

\begin{itemize}
\item \textbf{PV:} Physical Volume.
\item \textbf{VG:} Volume Group.
\item \textbf{LV:} Logical Volume.
\end{itemize}

Esta organización permite crecer, redistribuir y reorganizar almacenamiento de forma más flexible.
\subsubsection{Physical Volume}
\label{sec:orgcb1a241}

Un PV es un disco o partición preparado para LVM.

\begin{verbatim}
sudo pvcreate /dev/sdb
\end{verbatim}
\subsubsection{Volume Group}
\label{sec:org620baba}

Un VG es una reserva de almacenamiento formada por uno o varios PV.

\begin{verbatim}
sudo vgextend vg_datos /dev/sdb
\end{verbatim}

\texttt{vgextend} permite agregar un disco nuevo al grupo de volúmenes.
\subsubsection{Logical Volume}
\label{sec:org77173f8}

Un LV es el volumen que finalmente usa el sistema como si fuera una partición.

\begin{verbatim}
sudo lvextend -L +10G /dev/vg_datos/home
\end{verbatim}

Después de ampliar el LV, se debe ampliar el filesystem.
\subsubsection{resize2fs}
\label{sec:org07d276b}

En sistemas ext4, \texttt{resize2fs} permite que el filesystem use el nuevo espacio disponible.

\begin{verbatim}
sudo resize2fs /dev/vg_datos/home
\end{verbatim}

\texttt{lvextend} amplía el contenedor lógico y \texttt{resize2fs} amplía el contenido interno del sistema de archivos.
\subsubsection{Device mapper}
\label{sec:org74d9f69}

El device mapper es una capa del kernel que permite crear dispositivos virtuales a partir de dispositivos físicos.

Gracias a esta capa, LVM puede presentar volúmenes lógicos al sistema sin que las aplicaciones conozcan la distribución física real del almacenamiento.
\subsection{Conceptos de repaso}
\label{sec:org0a74afd}

\begin{itemize}
\item LVM desacopla almacenamiento físico y lógico.
\item \texttt{pvcreate} prepara discos.
\item \texttt{vgextend} amplía grupos de volúmenes.
\item \texttt{lvextend} amplía volúmenes lógicos.
\item \texttt{resize2fs} amplía el filesystem.
\item El filesystem no necesita conocer directamente el hardware físico.
\end{itemize}
\subsection{Máximas}
\label{sec:org0a48ca4}

\begin{itemize}
\item “El device mapper es el intermediario silencioso que hace posible que LVM funcione.”
\item “Extender un LV sin redimensionar el filesystem es ampliar el cuarto sin mover las paredes.”
\item “LVM permite que el almacenamiento crezca al ritmo del sistema.”
\item “Un volumen lógico no es solo espacio; es flexibilidad administrativa.”
\end{itemize}
\subsection{Sección libre}
\label{sec:orgc97843a}

Durante la práctica se relacionó LVM con instalaciones donde directorios como \texttt{/home}, \texttt{/var} y \texttt{/tmp} pueden configurarse como volúmenes lógicos.

Esto permite crecer partes específicas del sistema sin afectar todo el disco.

\newpage
\section{XAMPP, MySQL y archivos de base de datos}
\label{sec:orgb6362eb}

\subsection{Conceptos nuevos}
\label{sec:org6aab686}

\subsubsection{XAMPP}
\label{sec:org6724a34}

XAMPP integra herramientas necesarias para crear un entorno web local.

Incluye:

\begin{itemize}
\item Apache.
\item MariaDB o MySQL.
\item PHP.
\item Perl.
\end{itemize}

Sirve para probar aplicaciones web sin configurar cada servicio manualmente.
\subsubsection{Directorio de MySQL en XAMPP}
\label{sec:orga43d648}

En XAMPP, las bases de datos suelen estar en:

\begin{verbatim}
/opt/lampp/var/mysql
\end{verbatim}

Para revisar su contenido:

\begin{verbatim}
ls -la /opt/lampp/var/mysql
\end{verbatim}
\subsubsection{Archivos .ibd}
\label{sec:orgb20b496}

Los archivos \texttt{.ibd} pertenecen al motor InnoDB.

Almacenan:
\begin{itemize}
\item Datos reales de tablas.
\item Índices.
\item Información persistente del motor.
\end{itemize}

Esto muestra que una base de datos también depende del sistema de archivos.
\subsection{Conceptos de repaso}
\label{sec:orgd394e45}

\begin{itemize}
\item MySQL administra datos como archivos.
\item Linux controla permisos, rutas y almacenamiento.
\item \texttt{ls -la} muestra permisos y propietarios.
\item Los servicios web también dependen del filesystem.
\end{itemize}
\subsection{Máximas}
\label{sec:orgd6045d8}

\begin{itemize}
\item “Una base de datos también termina siendo archivo, permisos y almacenamiento.”
\item “XAMPP simplifica el laboratorio, pero no elimina la lógica real del servidor.”
\item “MySQL organiza datos; Linux decide cómo se guardan en disco.”
\end{itemize}
\subsection{Sección libre}
\label{sec:org10ed2c1}

Revisar el directorio de MySQL ayuda a conectar servicios con almacenamiento real. Aunque se interactúe con bases de datos mediante SQL, sus datos viven físicamente como archivos.

\newpage
\section{Realimentación y reflexión de aprendizajes}
\label{sec:org99a6992}

Esta semana permitió comprender cómo Linux automatiza servicios críticos mediante herramientas especializadas.

Chrony mostró que la sincronización de tiempo es indispensable para logs, auditorías, autenticación y coordinación entre equipos. Sin hora confiable, diagnosticar errores se vuelve más difícil.

AutoFS permitió entender una forma eficiente de montar recursos remotos. En lugar de montar todo durante el arranque, el sistema monta únicamente cuando alguien accede al recurso.

LVM reforzó la importancia de la abstracción en Linux. Separar PV, VG, LV y filesystem permite ampliar almacenamiento sin rediseñar todo el sistema.

La revisión de archivos internos de MySQL en XAMPP conectó los servicios con el sistema de archivos. Aunque MySQL trabaje a nivel de base de datos, sus datos dependen de archivos, permisos y almacenamiento físico.

En conjunto, la semana mostró que administrar Linux implica comprender la relación entre red, tiempo, almacenamiento, servicios y archivos.
\section{Comandos relevantes de la semana}
\label{sec:org05598ad}

\begin{verbatim}
systemctl status chronyd
sudo systemctl restart chronyd
chronyc sources -v
sudo firewall-cmd --add-service=ntp --permanent
sudo firewall-cmd --reload
sudo systemctl restart autofs
systemctl status autofs
sudo pvcreate /dev/sdb
sudo vgextend vg_datos /dev/sdb
sudo lvextend -L +10G /dev/vg_datos/home
sudo resize2fs /dev/vg_datos/home
ls -la /opt/lampp/var/mysql
\end{verbatim}
\section{Cierre}
\label{sec:org1579b22}

La Bitácora Semanal 11 integra tres ideas centrales: sincronización, automatización y abstracción.

Chrony representa la sincronización del sistema con una fuente confiable de tiempo. AutoFS representa la automatización del acceso a recursos remotos. LVM representa la abstracción del almacenamiento físico en estructuras lógicas más flexibles.

Estas herramientas demuestran que Linux permite construir sistemas más estables, escalables y administrables.
\end{document}
